% !TeX root = ../../Book2.tex
% 纲要标题：### 14.3 本原环与半本原环
% 纲要位置：计划纲要.md，第 1505 行。

\section{本原环与半本原环}
\label{b2:sec:1403}

模去 Jacobson 根以后，每个非零环元素至少能被某个简单模检测到。更强的问题是，能否用同一个简单模同时检测全部元素？本原环要求后者，半本原环只要求前者。它们与半单环的区别，将通过具体作用和无限乘积表现出来。

% 纲要要点（供目录核对）：
%
% * 忠实简单模与本原环；区分单个忠实作用、一族共同忠实作用及半单性
% * 半本原环的次直积结构；取极大左理想中最大的双边理想形成真正的本原商环
% * 半本原但非半单的实例；有限坐标满射与无限直积满射的区别
%

\subsection{忠实简单模与本原环}
\label{b2:subsec:1403:01}

\begin{definition}\label{b2:def:primitive-semiprimitive}
设 \(R\) 为非零含幺环。若 \(R\) 有忠实的幺性简单左模，则称它为\term[zuo-benyuan-huan]{左本原环}[left primitive ring]。若 \(J(R)=0\)，则称它为\term[banbenyuan-huan]{半本原环}[semiprimitive ring]。本节简称“本原”时均指左本原；右本原是另一条件，一般不能自动互换\footnote{右本原而非左本原的例子见\cite[pp. 473--475]{Bergman1964PrimitiveOneSided}；对该环取反环得到相反方向的例子。}。
\end{definition}

由\cref{b2:prop:radical-simple-annihilators}，半本原性等价于：对每个 \(0\ne r\in R\)，都存在一个幺性简单左模 \(S\) 和其中的向量 \(s\)，使 \(rs\ne0\)。这里 \(S\) 可以随 \(r\) 改变。本原性则要求先找到同一个简单模 \(S\)，使对每个 \(0\ne r\in R\)，都能在这个 \(S\) 中找到 \(s\) 满足 \(rs\ne0\)。因此每个本原环都半本原。对任意简单左模 \(S\)，商环 \(R/\operatorname{Ann}_R(S)\) 在 \(S\) 上的作用忠实且简单，所以它本原。

\ProseParagraphBreak
这里的“本原环”与下一章\cref{b2:def:primitive-idempotent}中的“本原幂等元”有不同的定义。前者要求一个忠实简单模，后者要求一个非零幂等元不能继续分成两个非零正交幂等元之和，不能仅凭同一个名称互换条件。

\begin{proposition}\label{b2:prop:commutative-primitive-field}
非零交换含幺环 \(R\) 左本原，当且仅当 \(R\) 是域。
\end{proposition}
\begin{proof}
若 \(S\) 忠实且简单，取 \(0\ne s\in S\)，满射 \(R\to S\)、\(r\mapsto rs\) 的核为极大理想 \(\mathfrak m\)，所以 \(S\cong R/\mathfrak m\)。交换性使这个商模的零化子恰为 \(\mathfrak m\)，忠实性便给出 \(\mathfrak m=0\)，故 \(R\) 是域。反向，一个域在自身上的正则模简单且忠实。
\end{proof}

在非交换情形，对任意除环 \(D\) 和整数 \(n\ge1\)，矩阵环 \(M_n(D)\) 都本原：列模 \(D^n\) 简单且忠实，简单性由矩阵单位移动非零坐标并乘除环元素得到。它在 \(n>1\) 时不是除环，说明上一个命题的交换性不能去掉。

\ProseParagraphBreak
非零半单环都半本原：左正则模是简单模的直和，而 \(J(R)\) 零化每个简单分量，所以 \(J(R)R=0\)，取右侧元素为 \(1\) 得 \(J(R)=0\)。但半单环未必本原：对任意域 \(k\)，环 \(k\times k\) 是两个域的乘积，因而半单；它是交换环却不是域，由上一个命题可知它不本原。

\subsection{半本原环的次直积结构}
\label{b2:subsec:1403:02}

\begin{definition}\label{b2:def:subdirect-product}
设 \(I\) 为指标集，\(R\) 与 \(R_i\)、\(i\in I\) 为含幺环，\(\phi:R\to\prod_{i\in I}R_i\) 为保幺环同态。若 \(\phi\) 单射，且它与每个因子投影的复合都满射，则称 \(\phi\) 把 \(R\) 表示为这些环的一个\term[cizhiji]{次直积}[subdirect product]。这里不要求 \(\phi\) 到整个直积满射。
\end{definition}

单射使像成为直积的一个子环，坐标满射使每个单独的因子都能取遍；不同坐标仍可能受到兼容条件限制。例如对任意域 \(k\)，对角嵌入 \(k\to k\times k\)、\(a\mapsto(a,a)\) 的两个坐标映射均满射，但它的像只包含两个坐标相等的元素。

\begin{theorem}\label{b2:thm:semiprimitive-subdirect}
非零含幺环 \(R\) 半本原，当且仅当它可以表示为一族左本原环的次直积。
\end{theorem}
\begin{proof}
对每个极大左理想 \(L\)，\(R/L\) 是简单左模，但 \(L\) 未必是双边理想，不能直接把它当作商环。为了得到一个在这个简单模上忠实作用的环商，应当模去作用的核，于是令
\[
I_L=\operatorname{Ann}_R(R/L)
=\{a\in R:aR\subseteq L\}.
\]
这个集合对加减封闭；若 \(aR\subseteq L\)，则对 \(r,x\in R\)，有 \((ra)x=r(ax)\in L\) 以及 \((ar)x=a(rx)\in L\)，故 \(I_L\) 是双边理想。取 \(x=1\) 可知 \(I_L\subseteq L\)。若双边理想 \(I\subseteq L\)，则每个 \(a\in I\) 都满足 \(aR\subseteq I\subseteq L\)，所以 \(I\subseteq I_L\)。这证明 \(I_L\) 是 \(L\) 中最大的双边理想。

商环 \(R/I_L\) 在 \(R/L\) 上的作用忠实，因为模去的恰是原作用的核；商前后子模相同，所以作用仍简单，故 \(R/I_L\) 本原。全部商映射合成
\[
\phi:R\longrightarrow\prod_L R/I_L.
\]
每个坐标映射满射，整体的核是 \(\bigcap_LI_L=J(R)\)，由\cref{b2:prop:radical-simple-annihilators}已知。故 \(J(R)=0\) 时它就是次直积表示。极大左理想组成集合，因此这里的乘积是集合编号的乘积。

反过来，设 \(\phi:R\to\prod_iR_i\) 为次直积表示，各 \(R_i\) 本原且半本原。每个坐标满射把 \(J(R)\) 映进 \(J(R_i)=0\)，这是\cref{b2:prop:radical-quotient}对其核所给的包含。于是 \(\phi\bigl(J(R)\bigr)=0\)，而 \(\phi\) 单射，故 \(J(R)=0\)。
\end{proof}

整数的模素数映射给出一个同时满足所有有限坐标要求、却不能任意指定无限坐标的例子：
\[
\mathbb Z\longrightarrow\prod_{p\text{ 素数}}\mathbb F_p.
\]
它单射，因为一个非零整数不可能被所有素数整除；各坐标显然满射。对任意有限多个不同素数，第一卷定理10.4.3的中国剩余定理允许同时任意指定这些坐标。但一个无限坐标族未必来自同一个整数：将奇素数坐标全取零、模二坐标取一就不可能提升。任何可能提升的整数须被所有奇素数整除，只能为零，又与模二条件矛盾。因此有限组同余条件的可解性，不能直接推广成对整个无限直积的满射性。

\subsection{非半单的半本原环}
\label{b2:subsec:1403:03}

\(\mathbb Z\) 的根为零，已在\secref{b2:sec:1401}由单位判据算出；它的正则模不半单，也已在\secref{b2:sec:1402}用不存在非零简单子模说明。另一方面，它不是域，由\cref{b2:prop:commutative-primitive-field}可知也不本原。所以这个半本原环既不本原，也不半单。

\ProseParagraphBreak
本原也不蕴含半单。令 \(k\) 为域，\(V\) 为有可数基 \(e_1,e_2,\ldots\) 的 \(k\) 向量空间，取 \(R=\operatorname{End}_k(V)\)。任意 \(0\ne v\in V\) 都可扩充为基；把 \(v\) 送到任意给定的 \(w\)，把这组基中的其余向量送到零，就唯一确定一个线性算子。无限维时扩充基使用第一卷附录A定理A.1.4的 Zorn 引理，和构造一般向量空间的基相同。故 \(Rv=V\)，自然模 \(V\) 简单；作用按定义忠实，所以 \(R\) 本原。

\ProseParagraphBreak
要检测这个环是否左 Artin，可以要求算子杀掉越来越多的指定基向量。这样的条件给出左理想
\[
L_n=\bigl\{\,T\in R:T(e_1)=\cdots=T(e_n)=0\,\bigr\}
\]
条件每增加一个，允许的算子就减少，形成无限严格降链 \(R\supsetneq L_1\supsetneq L_2\supsetneq\cdots\)。这些集合对加减和左乘稳定，因为复合 \(ST\) 仍将指定的基向量送到零；\(L_n\supsetneq L_{n+1}\) 的严格性可由只将 \(e_{n+1}\) 送到 \(e_1\)、将其余基向量送到零的算子核对，\(R\supsetneq L_1\) 则由恒等算子核对。因此 \(R\) 不是左 Artin 环，而\cref{b2:thm:semisimple-ring-characterizations}中的半单环均为左 Artin 环。这个例子说明本原性也不蕴含半单性，并且忠实简单模可能无限维。这里在有限组向量上指定算子的行为，正是下一节有限拓扑所记录的条件。
